\scanpage{23}
Il s'impose de regarder \unsure{églt}, si $A$ est ordinaire
\[
T_p(A)_{\text{et}}=\varprojlim\deleted{\unsure{T_p}}({}_{p^{\unsure{\nu|n}}}A)_{\text{et}}
\]
(NB. $T_p(A)$ est une extension
\diagram{\begin{tikzcd}[column sep=small]
0\arrow[r] & T_p(A)_{\mathrm{rad}}\arrow[r] & T_p(A)\arrow[r,dash] & T_p(A)_{\text{et}}\arrow[r,dash] & 0
\end{tikzcd}}
On sait, par dualité de Cartier, que dans la suite exacte correspondante
\diagram{\begin{tikzcd}[column sep=small]
\unsure{0}\arrow[r,dash] & T_p(A^*)_{\mathrm{rad}}\arrow[r,dash] & T_p(A^*)\arrow[r,dash] & T_p(A^*)_{\text{et}}\arrow[r,dash] & 0
\end{tikzcd}}
les termes médians sont duaux l'un de l'autre, et \deleted{plus précisément} cette dualité est compatible avec filtrations, donc
\[
\begin{aligned}
T_p(A)_{\mathrm{rad}}&\simeq D(T_p(A^*)_{\text{et}})\\
T_p(A^*)_{\mathrm{rad}}&\simeq D(T_p(A)_{\text{et}})
\end{aligned}
\]
donc la connaissance de $T_p(A)_{\text{et}}$ et $T_p(A^*)_{\text{et}}$ permet de reconstituer les \uline{deux} facteurs de $T_p(A)$, donc $T_p(A)$ si la base est parfaite [en général, il y aurait lieu d'expliciter la structure d'extension $\ldots$].

On aurait envie d'établir un isom can. \deleted{fonctoriel}
\mspage{22}
